\section{Introduction}
\label{sec:introduction}

A parcel-level crop map promises a catalog. The catalog is a decision: someone chose which crop
types the product will name and which it will not, and that choice is made long before any
accuracy figure is computed. In practice the choice is rarely reported as a choice. It arrives in a
sentence about how some classes had too few samples, and the evaluation that follows is computed
over whatever survived.

This matters because the metric that dominates the literature on imbalanced crop mapping is
macro-averaged F1, and macro-averaging divides by the number of classes. Removing a class the model
handled badly raises the macro average twice over: once by deleting a low per-class score from the
numerator, and once by shrinking the denominator. The second effect requires no model, no
mechanism, and no work. It is arithmetic. Yet a gain produced entirely by the second effect is
indistinguishable, in a results table, from a gain produced by a method.

We take that ambiguity as the object of study. We consider two mechanisms that a practitioner can
use to trade coverage for macro quality. The first shrinks the promised legend: the product names
$K$ of the $C$ available classes and declines to answer wherever its unrestricted prediction falls
outside them. The second keeps the full legend and abstains on the parcels where the predictor is
least confident, in the tradition of selective classification
\cite{chow1970reject,geifman2017selective}. Both reduce how much of the map is delivered; both
raise the macro figure;
and the question is how much of that rise either one is responsible for.

Our answer is a control, and it is almost embarrassingly cheap. Score the \emph{untouched}
predictor over the same shrunken legend while delivering everything. Whatever that control
recovers was never bought by the mechanism. On a first public benchmark of eighteen agronomic
classes, the control accounts for \num{88.3}\,\% of the apparent gain from shrinking the catalog
from eighteen classes to eight. On a second benchmark, in a different year, a different region and
with a different model family, it accounts for \num{95.1}\,\% at the pre-registered operating
point. And over the range where coverage remains complete, it accounts for \emph{all} of it: every
mechanism we test returns the same number, to four decimals, while the macro figure climbs by
\num{0.2447}.

Two further results follow from the same protocol. At matched coverage, abstaining on
low-confidence parcels outperforms shrinking the legend on both benchmarks, which is the opposite
of what the practice we set out to formalize assumes. And retiring classes by low support, which is
the criterion practitioners describe using, is the weakest of the three we test and is not monotone
in the number of classes retired: retiring more can make the delivered map worse.

The contributions are as follows. We formalize the two mechanisms and give a spatial-block protocol
that chooses the operating point outside the data used to measure it, scores both mechanisms over
the same class set so that their numbers are comparable, and never consults a ground-truth label to
decide which parcels to deliver. We quantify the decomposition of the apparent gain into
denominator and mechanism on two public benchmarks under an unchanged protocol. We order the three
retirement criteria with paired bootstrap intervals and a multiplicity correction. And we show that
a trained stacking meta-learner performs the same retirement implicitly and undeclared, and that
the difference between measuring it inside and outside the sample is enough to reverse the
conclusion about whether the ensemble helps at all.
